The code in Figure~\ref{table:fib} illustrates the use of this
framework with a simple Fibonacci computation. 

%% Add Fibonacci code in totality.

\begin{table}
{\footnotesize
\begin{verbatim}
public class FibG(n:Long) {
	static def fib(n:Long):Long=n<2? n: fib(n-1)+fib(n-2);
	
	class FibTQ implements TaskQueue{
		val bag = new ArrayListTaskBag[Long]();
		var result:Long=0;
		public def getTaskBag()=bag;
		public def getResult()= result;
		public def init(n:Long) {
			bag.bag().add(n);
		}
		public def process(var n:Long):Boolean {
			val b = bag.bag();
			for (var i:Long=0; bag.size() > 0 && i < n; i++) {
				val x = b.removeLast(); // constant time
				if (x < 2) result += x;
				else {
					b.add(x-1); // constant time
					b.add(x-2);
				}
			}
			return b.size()>0;
			
		}
		
	
		
		public def count() {
		
			return 0L;
		}
		
		public def merge(var _tb:glb.TaskBag):void {
			this.bag.merge(_tb as ArrayListTaskBag[Long]);
		}
		
	
		public def split():glb.TaskBag {
			return this.bag.split();
		}
		
		public def reduce():void {
			result= Team.WORLD.allreduce(result,Team.ADD); 
		}
		
		
	}
	//var finalResult:Long;
	public def run(n:Long):Long {
		
		val init = ()=>{ return new FibTQ(); };
		val glb = new GLB[FibTQ](init, GLBParameters.Default, true);
		
		Console.OUT.println("Starting...");
		var time:Long = System.nanoTime();
		val start = ()=>{ (glb.taskQueue()).init(n); };
		glb.run(start);
	
		 PlaceGroup.WORLD.broadcastFlat(()=>{
			(glb.taskQueue()).reduce();
		  
		});
		
		time = System.nanoTime() - time;
		Console.OUT.println("Finished.");
		Console.OUT.println("finalResult:" + glb.taskQueue().result);
		return glb.taskQueue().result;

	}
	public static def main(args:Rail[String]) {
		val N = args.size < 1 ? 10 : Long.parseLong(args(0));
		val result = new FibG(N).run(N);
		Console.OUT.println("fib(" + N + ")" + fib(N));
		Console.OUT.println("fib-glb(" + N + ")" + result);
		
	}
}
\end{verbatim}}
\caption{Fibonacci, with GLB}\label{table:fib}
\end{table}
